\documentclass[11pt,a4paper]{article}
\usepackage{../shared/luxcover}   % black cover (self-contained, TikZ logo, no assets)
\usepackage{../shared/proposal}   % \proposalhead helpers, stake environment, muted colour
\usepackage[utf8]{inputenc}\usepackage[T1]{fontenc}
\usepackage{amsmath,amssymb}
\usepackage{listings}\usepackage{xcolor}\input{../shared/lstlang}

\newcommand{\code}[1]{\texttt{\small #1}}
\newcommand{\prov}[1]{{\color{muted}\scriptsize\sffamily[#1]}}
\setlength{\emergencystretch}{3em}

\title{Bridge MOR, Start LMOR}
\author{Zach Kelling\\\textit{Lux Industries, Inc.}\\\texttt{research@lux.network}}
\date{2026}

\begin{document}
\luxcoverpage
\sloppy

\noindent{\color{muted}\footnotesize\sffamily ADOPT-13 \quad\textbullet\quad Hanzo AI \& Lux Industries \quad\textbullet\quad 2026}\par
\vspace{0.9ex}

\noindent\textbf{The number.}
Zero. MOR is a LayerZero OFT live on Ethereum, Arbitrum, and Base \prov{verified}, and its
collateral use on all three is zero \prov{verified}. It farms emissions or it trades. MOR is
\$1.95, market cap about \$17M, DEX depth \$1.4M with 88.8\% of it protocol-owned, trailing
24-hour volume \$8.9k \prov{verified}. A holder who wants cash sells. There is no way to borrow
against MOR without giving it up, and MOR is not on Lux.

\vspace{0.35ex}
\noindent\textbf{The problem, in their numbers.}
The Morpheus capital layer is sound: principal is withdrawable one to one and yield funds a real
buyback, about \$511K over the trailing year \prov{verified}. The gap is narrower. MOR's one live
use is farming its own emission, which runs about ten to one over real revenue \prov{verified}, and
the treasury's real yield is a single venue, stETH routed to Aave v3 through \code{DistributorV2.sol}
near 5\% \prov{verified}. So MOR has no non-speculative demand, its yield rides one counterparty, and
a holder cannot raise cash without selling into \$1.4M of mostly protocol-owned depth \prov{verified}.
None of this needs a Morpheus core change. It needs MOR reachable on a chain that lends against it.

\vspace{0.35ex}
\noindent\textbf{The fix, named to a repo.}
Two moves, near-term, no change to the Morpheus contracts. First, add MOR to \code{lux/bridge} and
list it on \code{lux/exchange} (lux.exchange), so MOR is bridgeable to Lux \prov{verified repo}. MOR
is already a LayerZero OFT \prov{verified}; extending an OFT to a new endpoint is a peer
configuration, not new token code \prov{derived}. Second, stand up LMOR: list MOR as collateral in
\code{lux/liquid}, the self-repaying lending engine. Deposit MOR into \code{VaultV2}, mint a
transferable synthetic against it up to the position's floor, and the collateral's yield is harvested
and burned against the debt, so the balance owed falls on its own \prov{verified repo}.
\code{lux/liquid} routes that yield through about a dozen risk-capped adapters, among them Aave,
Compound, Lido, EtherFi, Pendle, Morpho, Yearn, Ethena, EigenLayer, and Maker DSR
\prov{verified repo}, not Aave alone.

\vspace{0.35ex}
\noindent\textbf{The outcome, projected.}
A MOR holder borrows against MOR without selling, the loan self-repays from yield, and that yield is
spread across a dozen capped adapters instead of one near 5\% \prov{verified}. Assumptions stated.
Size an initial LMOR cap to what MOR depth can liquidate, near \$600k of quote-side depth today
\prov{verified}: at a conservative 33\% floor for a thin, volatile collateral \prov{illustrative},
that is about \$0.2M of spendable liquidity minted against MOR that produces \$0 spendable for its
holder today \prov{derived}, and the cap grows as the book does. MOR gains its first non-speculative
demand, sourced from lending rather than from the emission schedule. This starts against the existing
Arbitrum and Base MOR liquidity the moment MOR bridges: Lux gears up the bridgeable token and the
LMOR market on its own side and shows a working market before any capital migrates.

\vspace{0.35ex}
\noindent\textbf{Relevance, and first step.}
This is the fastest path to a real use for MOR, and the path is Lux infrastructure end to end, the
bridge, the exchange, and the lending engine, not a Morpheus rebuild. It needs Morpheus permission
for nothing but one peer configuration \prov{verified}. First step: Lux lists the MOR pair on
\code{lux/dex} (LX), stands up the LMOR market in \code{lux/liquid} against existing Arbitrum and Base
liquidity, and publishes depth, the mint, and the self-repay against live figures; Morpheus adds the
Lux endpoint as an OFT peer when it chooses. The lux.exchange offering venue depends on an SEC
approval Lux is pursuing; it is pending, not granted \prov{pending}, and nothing here is a
solicitation or an offer of a security.

\vspace{0.35ex}
\noindent\textbf{What this does not solve.}
Bare MOR is borrowable, but a MOR position only retires itself if MOR is put to work in a yield
strategy of its own; self-repayment is a property of yield-bearing collateral, not of the listing
\prov{verified}. The synthetic holds its peg through the redemption schedule and collateral quality,
not by decree. The projected liquidity is bounded by real MOR depth, thin today, and grows only as
the book does.

\vspace{0.5ex}
\noindent{\color{muted}\footnotesize\sffamily
ADOPT-13. Rails: \code{lux/bridge}, \code{lux/exchange} at lux.exchange, \code{lux/liquid}
(\code{VaultV2} and a dozen adapters). MOR is a LayerZero OFT on Ethereum, Arbitrum, and Base. Pairs
with ADOPT-04 (lending) and ADOPT-06 (order-book venue). Every citation checked against the working
tree.\quad\textbullet\quad \texttt{research@lux.network}}

\end{document}
